\section{Reduction of the numerical assertion to algebraic curves}
\label{sec:algebraic-reduction}

We now use the Wahl invariant $\beta=-P_{\mathrm{loc}}^2$ defined above.
The objective of this section is to transfer constancy of this scalar
invariant from algebraic curve families to arbitrary holomorphic disk
families.  The individual right-equivalences of
Lemma~\ref{lem:individual-jets} are sufficient; no simultaneous analytic
equivalence of the original family with a polynomial family is asserted.

Fix $\mu>0$ and $N\geq\mu+1$.  Let $V_N$ be the affine space of complex
polynomials in three variables of degree at most $N$, with zero constant
and linear terms, and put
\[
 U_\mu=\{g\in V_N:\dim_{\mathbb C}\mathbb C\{x\}/J(g)=\mu\}.
\]

\begin{lemma}[The precise jet locus]
\label{lem:mu-jet-locus}
The subset $U_\mu$ is locally closed algebraic.  On $U_\mu$, the function
$g\mapsto\beta(V(g),0)$ has finite image and is algebraically
constructible.
\end{lemma}

\begin{proof}
Work first in the finite-dimensional algebra
$A_\mu=\mathbb C[x_1,x_2,x_3]/\mathfrak m^{\mu+1}$.  The ideal generated
by the partials of $g$ is the image of a linear map
$A_\mu^{\oplus3}\to A_\mu$ whose matrix entries depend polynomially on
the coefficients of $g$.  The condition that its cokernel have
dimension $\mu$ is a locally closed rank condition.

This condition tests the full local length, including exclusion of
infinite length.  Indeed, for the local algebra
$A=\mathbb C\{x\}/J(g)$ with maximal ideal $\mathfrak n$, suppose
$\dim_{\mathbb C}A/\mathfrak n^{\mu+1}=\mu$.  The $\mu+1$ successive
layers $\mathfrak n^j/\mathfrak n^{j+1}$, $0\leq j\leq\mu$, cannot
all be nonzero.  If a layer vanishes, Nakayama's lemma applied to the
finitely generated ideal $\mathfrak n^j$ shows that this ideal is zero.
Consequently $A$ already has length $\mu$.  The converse follows from
the same filtration.  Polynomial, localized algebraic, and convergent
power-series computations give the same truncated quotient.  This
proves the first assertion.

For constructibility, let $B$ be any irreducible locally closed
subvariety of $U_\mu$; replace it first by a smooth dense open.
Consider its universal hypersurface with the origin section.  After
another dense open restriction on $B$, there is a Zariski neighborhood
of this section on which the relative singular locus is exactly the
section.  To see this, isolatedness in each fiber implies that the
section is an irreducible component of the relative critical locus:
a larger component containing it would have a positive-dimensional
fiber at a general section point, by the dimension inequality.
The other components meet the section over proper closed subsets of
$B$.  Remove those subsets from $B$, then remove those components
from the total space.  The resulting total hypersurface is
Cohen--Macaulay and regular off a subset of codimension two, hence
normal.  Its fibers have dimension two, so miracle flatness over the
smooth base gives flatness; it is smooth over $B$ off the section.

Take a projective resolution, isomorphic off the section, which
principalizes its ideal and has simple normal crossing reduced
exceptional divisor $E$.  Characteristic-zero resolution of
singularities and principalization provide this construction; see
\cite[Theorems~1.0.1--1.0.3]{Wlodarczyk2005}.
Generic smoothness, applied to the resolution space and every
intersection of components of $E$, permits a further dense open
restriction on $B$ such that the resolution is smooth over $B$ and
$E$ is a relative simple normal crossing divisor.  Nondominating
exceptional strata are discarded in making this restriction.
Since the resolution is projective and all of $E$ lies over the
section, $E$ and all intersections of its components are proper over
$B$.

The fibers are now good resolutions of the surface germs at the
origins.  Their weighted exceptional dual graphs are locally constant
in the complex topology, up to permutation of the vertices.
In fact smooth proper variation preserves the connected components
and genera of the exceptional curves, as well as their intersection
points and incidence.  Their self-intersections are degrees of the
normal line bundles obtained by restricting the divisor line bundles
on the resolution space.  These degrees are locally constant in
smooth proper curve families.  If a total divisor has several
components in a fiber, they are disjoint because that relative divisor
is smooth, so the same degree calculation applies to each component.

The weighted graph determines $\beta$.  It determines the exceptional
intersection matrix, and adjunction gives
$K\cdot E_i=2g(E_i)-2-E_i^2$ for each exceptional curve.  Thus it
determines all intersections of $K+E$ with the exceptional curves,
and hence its relative Zariski decomposition and the square of its
numerical exceptional representative.  An irreducible smooth complex
base is connected, so $\beta$ is constant on this dense open of $B$.
Noetherian induction on the remaining closed subsets yields a finite
stratification of $U_\mu$ on which $\beta$ is constant.
\end{proof}

\begin{proposition}[Scalar curve reduction]
\label{prop:scalar-curve-reduction}
Assume the following algebraic assertion: for every polynomial family
in three variables with regular coefficients on a smooth complex
algebraic curve $C$, with zero constant and linear terms and constant
positive Milnor number at the origin, the Wahl invariant at a marked
non-log-canonical fiber equals its value on the nearby general fibers.
Then $\beta(X_t,0)$ is constant near $t=0$ in every holomorphic
$\mu$-constant family whose central surface germ is non-log-canonical.
\end{proposition}

\begin{proof}
Take $N\geq\mu+1$ and consider the holomorphic coefficient map
$a:\Delta\to V_N$, $a(t)=j^Nf_t$.  By
Lemma~\ref{lem:individual-jets}, it takes values in $U_\mu$ and
preserves the analytic type of every individual surface germ, hence
its Wahl invariant.  Write $a_0=a(0)$ and $\beta_0=\beta(V(a_0),0)$.

If constancy fails on every smaller disk, $a_0$ lies in the closure
of the constructible set
\[
 Z_{\ne}=\{g\in U_\mu:\beta(V(g),0)\ne\beta_0\}.
\]
The closure here may be taken in the Zariski topology: an accumulating
sequence of coefficient vectors already forces this inclusion.
Choose an irreducible locally closed subset of $Z_{\ne}$ whose closure
in $U_\mu$ contains $a_0$.  After restricting to an affine neighborhood
of $a_0$, its closure is an irreducible affine variety $Z$ of positive
dimension, containing a dense open $Z^\circ\subset Z_{\ne}$.

There is an algebraic curve in $Z$ through $a_0$ whose general point
lies in $Z^\circ$.  For completeness, choose a finite Noether
normalization $\pi:Z\to\mathbb A^d$ and translate so that
$\pi(a_0)=0$.  The image of $Z\setminus Z^\circ$ is a proper closed
subset.  Choose a line through $0$ not contained in that image.
Its inverse image is cut out in $Z$ by $d-1$ equations, so every
irreducible component through $a_0$ has dimension at least one.
Finiteness of $\pi$ gives the opposite inequality and prevents a
one-dimensional component from mapping to a point.  Such a component
therefore dominates the line and has general point in $Z^\circ$.
Normalize it, choose a point over $a_0$, and remove the finitely many
other points needed to obtain the desired marked smooth curve $C$.
The inclusion into coefficient space supplies regular polynomial
coefficients on $C$.

Every member of this curve lies in $U_\mu$, its marked fiber is the
polynomial $a_0$ with non-log-canonical surface germ, and its general
Wahl invariant differs from $\beta_0$.  This contradicts the assumed
algebraic assertion.  Therefore the invariant is constant in the
original analytic family after shrinking.
\end{proof}

Appendix~\ref{app:local-numerical} gives a second curve comparison,
using irreducible components of $U_\mu$ directly; it does not require
constructibility of the scalar invariant.
